\cotefile{150}
\scanpage{85}

\noindent\uline{16}\quad $\Sigma_J$ et $\Sigma_{J,L}$.

\margin{Un point de vue peut-être plus commode pour définir les $\Sigma_i$, $\Sigma_{i_0\ldots i_n}$, $\Sigma_J$, $\Sigma_J^{L}$ est celui des ``éclatés orientés'', mieux adapté à des formules du type $X=\underset{d\in\mathrm{Drap}(I)}{\uline{\lim}}\Sigma_d$. On prend des $\Sigma_i\supset X_i^*$ et non $\deleted{\unsure{\Sigma_i}}\subset X_i^*$, \unsure{idem pour} $\Sigma_J$, $X_J^*=\Sigma_J\setminus\dot\Sigma_J$\ldots}

Soit $J$ une bande de $I$, on définit
\[
(70)\qquad
\begin{array}{@{}l@{}}
\Sigma_J\subset X_J^*\\[-.4ex]
\hspace{.6em}\Vert\\[-.4ex]
X_J^*\setminus X_J^*\cap\displaystyle\bigcup_{i\in\bar J\setminus J}\mathring T_i
\end{array}
\]
\[
\begin{array}{@{}c@{\quad}l@{}}
\deleted{\Vert}&\Vert\quad\text{\uline{NB}}\\[-.4ex]
\deleted{\displaystyle\bigcup_{j\in J}\Sigma_j}&
X_J^{\deleted{\unsure{*}}}\setminus X_J^{\deleted{\unsure{*}}}\cap\displaystyle\bigcup_{i\in\bar J\setminus J}T_i
\end{array}
\]
\margin{NB $\bar J=$ idéal engendré par $J$ $=\{i\in I\mid\exists j\in J\ \text{avec}\ i\leq j\}$}
\margin{\uline{NB} On a posé $X_J=\bigcup_{j\in J}X_j\ (=\bigcup_{j\in\bar J}X_j^*=X_{\bar J}^*)$}

On a donc
\[
(71)\quad\left\{
\begin{aligned}
\Sigma_{\phi}&=\phi,\quad\Sigma_{\{j\}}=\Sigma_j,\quad\Sigma_I=X\\
\Sigma_J&=X_J\quad\text{si $J$ un idéal i.e. $J=\bar J$}
\end{aligned}\right.
\]

On a même, sauf erreur
\[
(72)\qquad\Sigma_J\simeq\underset{i,j\in J}{\lim_{\longrightarrow}}\left(\Sigma_j,\Sigma_{ij},\deleted{\unsure{b}}(s_{ij},b_{ij}:\right.
\]
\diagram{\begin{tikzcd}[ampersand replacement=\&,column sep=small,row sep=small]
\&\Sigma_i\\
\Sigma_{ij}\arrow[ur]\arrow[dr]\&\\
\&\Sigma_j
\end{tikzcd}}
\[
\left.\vphantom{\Sigma_{ij}}\right)\qquad\simeq\underset{d\in\mathrm{Drap}(J)}{\uline{\lim}}\Sigma_d
\]
\margin{\uline{NB} $\Sigma_J$ est muni d'une stratification de type $J$, que \unsure{l'on déduit} des $(\Sigma_j,\Sigma_{jh},s_{jh},b_{jh})$ avec $j,h\in I$}
\deleted{Par suite \unsure{d'un}}
\[
(73)\qquad\Sigma_J\overset{\overset{h}{\sim}}{\hookrightarrow}X_J^*\qquad(\text{équivalence d'homotopie})
\]

On pose
\[
(74)\qquad\begin{aligned}
\dot\Sigma_J&=X_J\cap\left(\text{bord de }\underbrace{T_{\bar J\setminus J}}_{\displaystyle\bigcup_{i\in\bar J\setminus J}T_i}\right)\\
&=\bigcup_{\substack{i<j\\i,j\in J}}\operatorname{Im}(\Sigma_{ij}\to\Sigma_J)\\
&=\underset{\substack{d\in\mathrm{Drap},\ d\text{ \unsure{ayant|avec} au moins 2 termes}}}{\lim_{\twoheadrightarrow}}\Sigma_d
\end{aligned}
\]
% Thin lower foliation, in another hand; not the author's page number.
\margin{84}
